% =============================================================================
% Experimental Setup — Whisper Hallucination paper (Interspeech 2026)
% Paste into your Overleaf project's methodology/experiments section.
% =============================================================================

\section{Experimental Setup}
\label{sec:experiments}

We investigate how different patterns and quantities of label noise during
fine-tuning cause distinct failure modes in Whisper Large~v3.
Our study comprises two complementary experiments: a \emph{noise-pattern}
experiment that compares four structurally different mismatch types at a fixed
contamination rate, and a \emph{dose--response} sweep that varies the
contamination rate to identify the onset threshold of each failure mode.

% ---------------------------------------------------------------------------
\subsection{Data and Model}
\label{sec:data}

All experiments share a common setup.
We draw a 120\,k-utterance subset (stratified random sample, seed\,=\,42)
from Common Voice English~22~\cite{commonvoice}, totalling approximately
360 hours of read speech.
The same subset serves as the clean base for every condition, ensuring that
differences across runs are attributable only to the injected noise.
Development and test sets are the standard Common Voice splits (16.4\,k and
16.4\,k utterances, respectively).

We fine-tune \textbf{Whisper Large~v3}~\cite{radford2023robust} using
LoRA~\cite{hu2022lora} with rank $r{=}16$, $\alpha{=}32$,
dropout\,=\,0.1, applied to all attention projections and feed-forward
layers (\texttt{q\_proj}, \texttt{k\_proj}, \texttt{v\_proj},
\texttt{o\_proj}, \texttt{fc1}, \texttt{fc2}).
All models train for 5 epochs with AdamW (lr\,=\,$10^{-4}$,
500 warmup steps), per-device batch size 16, gradient accumulation 2
(effective batch 32), fp16 mixed precision, and gradient checkpointing.
The best checkpoint is selected by validation WER.

% ---------------------------------------------------------------------------
\subsection{Experiment~1: Noise-Pattern Comparison}
\label{sec:exp1}

We compare four structurally distinct label-noise patterns, each injected
at a fixed 8\% contamination rate (${\sim}$9\,600 noisy utterances added
to the 120\,k clean set), plus a clean baseline.
Table~\ref{tab:noise-patterns} summarises the configurations.

\begin{table}[t]
\centering
\caption{%
  Noise-pattern configurations for Experiment~1.
  Each adds ${\sim}$9\,600 mismatched utterances to the 120\,k clean training set.
  \emph{Audio} and \emph{Text} columns describe the diversity of the noisy
  portion only.
}
\label{tab:noise-patterns}
\renewcommand{\arraystretch}{1.15}
\begin{tabular}{@{}lllp{4.0cm}@{}}
\toprule
\textbf{Config} & \textbf{Audio} & \textbf{Text} & \textbf{Description} \\
\midrule
Base & \multicolumn{3}{l}{Clean labels only (control)} \\
UU   & Unique   & Unique   & Each noisy sample has a distinct audio
                              paired with a distinct unrelated sentence. \\
RR   & Repeated & Repeated & 10 audio--text mismatch pairs, each
                              repeated ${\sim}$960\,\texttimes. \\
RU   & Repeated & Unique   & 10 audio clips, each time with a
                              different unrelated sentence. \\
UR   & Unique   & Repeated & Unique audio clips, all paired with one
                              of 10 repeated sentences. \\
\bottomrule
\end{tabular}
\end{table}

The four patterns probe orthogonal failure mechanisms:
\begin{itemize}[nosep]
  \item \textbf{UU} tests whether the model learns to \emph{hallucinate}
        plausible but incorrect text when every noisy sample is unique.
  \item \textbf{RR} tests whether repeated exposure to the same audio--text
        pair causes the decoder to memorise and \emph{repeat} fixed outputs.
  \item \textbf{RU} isolates the effect of \emph{audio repetition}: the same
        waveform appears with many different labels, potentially destabilising
        the encoder representation.
  \item \textbf{UR} isolates \emph{text repetition}: the decoder sees the same
        sentence paired with varied acoustics, potentially biasing it toward
        high-frequency outputs.
\end{itemize}

% ---------------------------------------------------------------------------
\subsection{Experiment~2: Dose--Response Sweep}
\label{sec:exp2}

Experiment~1 uses a single, arguably arbitrary contamination rate.
To establish \emph{how much} noise is needed to elicit each failure mode,
Experiment~2 sweeps the noise ratio across
$\{1, 2, 5, 10, 20, 50\}\%$
for the two most contrasting patterns: \textbf{UU} (diverse mismatches
$\to$ hallucinations) and \textbf{RR} (repeated mismatches $\to$
repetition loops).
This yields 12 training runs plus the shared 0\% baseline, for a total
of 13 data points per metric.
All other hyperparameters are identical to Experiment~1.

The resulting dose--response curves (Section~\ref{sec:results}) reveal
(i) whether failure modes emerge gradually or with a sharp onset,
(ii) the minimum contamination level that produces measurable degradation,
and (iii) whether the two failure modes saturate at different rates.

% ---------------------------------------------------------------------------
\subsection{Evaluation Framework}
\label{sec:eval}

Each trained model is evaluated on the clean Common Voice test set using
three complementary metrics:

\paragraph{Word accuracy (WAcc).}
$\mathrm{WAcc} = 1 - \mathrm{WER}$, computed with the \texttt{jiwer}
library after standard text normalisation (lower-casing, punctuation removal).
This measures \emph{lexical correctness} with respect to the ground-truth
transcription.

\paragraph{Sequence plausibility.}
We score each hypothesis with Qwen3-1.7B~\cite{qwen3}, computing the
mean per-token log-probability of the generated sentence.
A plausible hallucination will have \emph{high} plausibility despite
\emph{low} WAcc, separating it from random decoder noise.
Using an open-source LM ensures full reproducibility.

\paragraph{$N$-gram repetition rate.}
For each hypothesis we count repeated 3-grams (trigrams appearing more
than once) and report the fraction of test utterances containing at least
one repeated trigram.
This captures the \emph{looping} behaviour that characterises
repetition-type failures.

\paragraph{Input perturbation.}
To stress-test robustness, we additionally evaluate each model under two
acoustic perturbations applied to the test audio:
(a)~\emph{onset noise}: additive Gaussian noise ($\sigma{=}0.1$) injected
into the first 0.5\,s; and
(b)~\emph{full noise}: additive noise applied to the entire waveform.
Perturbation-induced degradation quantifies how aggressively each failure
mode is triggered by out-of-distribution inputs.

\paragraph{Summary.}
Table~\ref{tab:experiment-summary} provides an overview of both experiments.

\begin{table}[t]
\centering
\caption{Overview of the two experiments.}
\label{tab:experiment-summary}
\renewcommand{\arraystretch}{1.15}
\begin{tabular}{@{}lcc@{}}
\toprule
                          & \textbf{Exp.\,1} & \textbf{Exp.\,2} \\
\midrule
Noise patterns            & UU, RR, RU, UR   & UU, RR            \\
Noise ratios              & 8\%              & 1--50\%           \\
Training runs             & 5                & 13                \\
Clean training utterances & 120\,k           & 120\,k            \\
Dev / Test utterances     & 16.4\,k / 16.4\,k & 16.4\,k / 16.4\,k \\
Model                     & \multicolumn{2}{c}{Whisper Large v3 + LoRA} \\
Epochs                    & \multicolumn{2}{c}{5}               \\
Effective batch size       & \multicolumn{2}{c}{32}              \\
Metrics                   & \multicolumn{2}{c}{WAcc, plausibility, repetition, perturbation} \\
\bottomrule
\end{tabular}
\end{table}
